% \documentclass[sigconf]{acmart}

\documentclass[nonacm,sigconf]{acmart}

\usepackage{paralist}
\usepackage{microtype}

% Tighter line-breaking inside compactitem and inline \texttt{}.
\sloppy
\emergencystretch=2em



%%
%% end of the preamble, start of the body of the document source.
\begin{document}

%%
%% The "title" command has an optional parameter,
%% allowing the author to define a "short title" to be used in page headers.
\title{Time-Delayed Publicly Verifiable Quantum Computation with Classical Verifiers}

%%
%% The "author" command and its associated commands are used to define
%% the authors and their affiliations.
\author{Ameer Mohammed}
\email{ameer.mohammed@ku.edu.kw}
\affiliation{%
  \institution{Kuwait University}
  \city{Kuwait City}
  \country{Kuwait}}

\author{Aydin Abadi}
\email{aydin.abadi@ncl.ac.uk}
\affiliation{%
  \institution{Newcastle University}
  \city{Newcastle upon Tyne}
  \country{United Kingdom}}

\author{Jaffer Mahdi}
\email{jaffer.mahdi@ku.edu.kw}
\affiliation{%
  \institution{Kuwait Oil Company}
  \city{Kuwait City}
  \country{Kuwait}}
%%
%% This command processes the author and affiliation and title
%% information and builds the first part of the formatted document.
\maketitle


%%%%%%%%%%%%%%%%%%%%%%%%%%%%%%%%%%%%%%%%%%%%%%%%%%%%%%%%%%%%%%%%%%%%%
% Artifact Appendix for ACM CCS 2026.
% Please limit your Artifact Appendix to no more than two pages.
%%%%%%%%%%%%%%%%%%%%%%%%%%%%%%%%%%%%%%%%%%%%%%%%%%%%%%%%%%%%%%%%%%%%%

%% This template uses list-making environments from the `paralist`
%% LaTeX package.  Add the following to the preamble of your LaTeX
%% document (somewhere before `\begin{document}`:
%%
%% \usepackage{paralist}

%%%%%%%%%%%%%%%%%%%%%%%%%%%%%%%%%%%%%%%%%%%%%%%%%%%%%%%%%%%%%%%%%%%%%

\appendix


\section{Artifact Appendix}

This appendix supports the paper's central claim: for the linear
systems studied, Enhanced Hybrid HHL finishes before the time-lock
puzzle (TLP) protecting the verification key can be solved.  The
artifact packages three self-contained Python experiments behind a
single entry point and is intended to be reproducible on any basic computer.

\subsection{Abstract}

A single Python project at the repository root.  The entry
point \texttt{run\_experiments.py} selects among three experiments
via CLI flags (\texttt{--hhl-only}, \texttt{--tlp-only},
\texttt{--scaling-only}) and writes PNG plots to
\texttt{visualization/outputs/} plus JSON data files to the working
directory.  All inputs (linear system, sweep ranges,
clock-qubit count) live in the \texttt{INPUTS} block at the top of the file
. Licensed under Creative Commons Attribution 4.0
International (CC-BY~4.0); see the Zenodo record for the
deposited copy. 
\texttt{run\_experiments.py}.

\subsection{Description \& Requirements}

\subsubsection*{Artifact layout}

\begin{compactitem}
\item \texttt{README.md} --- overview and CLI reference.
\item \texttt{run\_experiments.py} --- entry point and \texttt{INPUTS} block.
\item \texttt{enhanced\_hybrid\_hhl/} --- Enhanced Hybrid HHL solver.
\item \texttt{timelock\_puzzle/} --- TLP garbling, generation, solving.
\item \texttt{crgc\_garbled\_circuits/} --- reusable garbled-circuit primitives.
\item \texttt{visualization/} --- plotting code and PNG outputs.
\item \texttt{circuits/} --- Bristol-format SHA-256 used inside the TLP.
\item \texttt{requirements.txt} --- Python dependencies.
\end{compactitem}

\subsubsection{Security, privacy, and ethical concerns}

The artifact performs only local, numerical simulations
on a Qiskit \texttt{AerSimulator} backend.  It makes no network calls,
does not modify the host system outside the working directory and
\texttt{visualization/outputs/}. It also does not do anything that might be 
considered unethical.

\subsubsection{How to access}

The artifact is permanently archived at Zenodo under DOI
\texttt{10.5281/zenodo.20751485} (record
\href{https://zenodo.org/records/20751485}{zenodo.org/records/20751485}),
licensed under Creative Commons Attribution 4.0 International
(CC-BY~4.0).  Download the \texttt{Quantum\_vertification\_artifact.zip}
from the record to obtain the full repository.

\subsubsection{Hardware dependencies}

A regular laptop suffices.  Reference machine: 2024 Macbook Pro, M3 Max,
36\,GB RAM.  No GPU, FPGA, QPU, or cloud account is required; all
quantum circuits run on the local Qiskit \texttt{AerSimulator}.
Absolute times will vary with the host CPU, but the qualitative
relationships between HHL, TLP, and the paper's claims are
hardware-independent.

\subsubsection{Software dependencies}

Python $\geq 3.10$ on macOS, Linux, or Windows.  The five direct
dependencies are pinned in \texttt{requirements.txt}:
\texttt{qiskit\,$\geq$\,1.0.0}, \texttt{qiskit-aer\,$\geq$\,0.14.0},
\texttt{numpy\,$\geq$\,1.26.0},
\texttt{matplotlib\,$\geq$\,3.8.0}, \texttt{scipy\,$\geq$\,1.12.0}.
The paper's timings were collected on AWS Braket SV1; this artifact
uses Qiskit \texttt{AerSimulator} locally.  Absolute timings differ
across simulators, but the relative ordering is preserved.
\textbf{Not validated:} the paper's exact Braket SV1 wall-clock
numbers (they require paid cloud access); this artifact validates
the qualitative relationships, not the absolute Braket times.

\subsubsection{Benchmarks}

None required.  Toy linear systems are used for the experiments; the TLP's sequential function uses the SHA-256
Bristol-format circuit shipped at
\texttt{circuits/sha256.txt}.

\subsection{Set Up}

\subsubsection{Installation}

Using the \texttt{uv} package manager (project-preferred):
\begin{verbatim}
uv sync
\end{verbatim}
Equivalently, with \texttt{pip}:
\begin{verbatim}
pip install -r requirements.txt
\end{verbatim}
This resolves all dependencies.

\subsubsection{Basic test}

\begin{verbatim}
uv run python run_experiments.py --hhl-only
\end{verbatim}
On success the script prints an HHL total time on the order of
seconds and writes \texttt{visualization/outputs/hhl\_results.png}.
If the file appears and the graph shows both hhl matrix components, the artifact
is ready for the full evaluation below.

\subsection{Evaluation Workflow}

\subsubsection{Major claims}

\begin{compactitem}
\item[(C1):] \emph{Enhanced Hybrid HHL on the 2$\times$2 system
  $A=[[2,1],[0,1.5]]$, $b=[1,2]$ produces a sign-agnostic (negative values
  not captured by result) solution with finite L2 error and values close to the classical
  linear system of eqns solution, completing in a few seconds on the local Qiskit
  \texttt{AerSimulator}.}
  Proven by (E1); see
  \texttt{visualization/}\-\texttt{outputs/}\-\texttt{hhl\_results.png}.

\item[(C2):] \emph{TLP \texttt{PSolve} time scales linearly in $T$
  (intrinsic to the protocol: a garbled SHA-256 is evaluated
  sequentially $T$ times).}
  Proven by (E2); see \texttt{tlp\_benchmark\_plot.png} and
  \texttt{tlp\_benchmark\_results.json::reveal\_times}.

\item[(C3):] \emph{For $N \in \{2, 4, 8\}$ with calibrated
  $T(N) \in \{250, 800, 4000\}$, the projected TLP \texttt{PSolve}
  time exceeds the measured HHL average time at every $N$ (the
  paper's main quantitative claim).  Absolute times vary with the
  host CPU; the TLP $>$ HHL ordering is hardware-independent. As it 
  can be selected and configured by the user, by changing the TLP time parameter $T$ for a given
  setup}
  Proven by (E3); see \texttt{scaling\_results.png} and
  \texttt{scaling\_results.json::tlp\_solve\_times} vs.\
  \texttt{hhl\_avg\_times}.
\end{compactitem}

\subsubsection{Experiments}

\begin{compactitem}
\item[(E1):] \emph{[Enhanced Hybrid HHL, $<$1~min + $<$1~min + $<$1\,MB output]}.
  Solves the 2$\times$2 system $A=[[2,1],[0,1.5]]$, $b=[1,2]$ (set
  by \texttt{HHL\_A}, \texttt{HHL\_B}) using Enhanced Hybrid HHL on
  the local \texttt{AerSimulator}.

  \begin{asparadesc}
  \item[Preparation:] \texttt{uv sync} (or
    \texttt{pip install -r requirements.txt}).

  \item[Execution:]
    \texttt{uv run python run\_experiments.py --hhl-only}.

  \item[Results:] Inspect
    \texttt{visualization/}\-\texttt{outputs/}\-\texttt{hhl\_results.png}
    and the printed \texttt{l2\_error}, \texttt{circuit\_depth},
    \texttt{total\_time}.  Pass criterion: Each component's value is 
    similar to the classical solution, and the \texttt{l2\_error} is
    a small finite value (a small apparent discrepancy is expected since HHL
    recovers only amplitude magnitudes, not signs).
  \end{asparadesc}

\item[(E2):] \emph{[TLP benchmark, $<$1min + $<$2min + $<$1\,MB output]}.
  Times the three TLP phases (garbling, generation, solving) for
  $T\in\{1,2,3,4\}$ with 3 trials per $T$ (set by
  \texttt{TLP\_T\_MIN}, \texttt{TLP\_T\_MAX}, \texttt{TLP\_TRIALS}).
  Defaults finish in a few minutes on the reference machine.

  \begin{asparadesc}
  \item[Preparation:] Same as E1.

  \item[Execution:]
    \texttt{uv run python run\_experiments.py --tlp-only}.

  \item[Results:] Inspect \texttt{tlp\_benchmark\_plot.png} and
    \texttt{tlp\_benchmark\_results.json}.  Pass criteria: all
    trials pass \texttt{recovered == secret\_bit}, and
    \texttt{reveal\_times} grows linearly with $T$ with a small intercept and error.
  \end{asparadesc}

\item[(E3):] \emph{[Scaling, $<$1~min + $<$20~min + $<$1\,MB output]}.
  For each $N\in\{2,4,8\}$ (see \texttt{SCALING\_N\_VALUES}),
  generates 10 random $N{\times}N$ matrices, runs Enhanced Hybrid
  HHL on each, calibrates TLP \texttt{PSolve} at
  $T=1,\ldots,5$ on this machine, fits a linear model, and projects
  \texttt{PSolve} to the calibrated $T(N)\in\{250,800,4000\}$ (see
  \texttt{N\_TO\_T}).  Pass \texttt{--tlp-actual} to run the
  full pipeline at the large $T(N)$ (takes a lot of time depending on T; only needed
  to verify the linear projection), else you can keep it as is to allow for the 
  program to estimate the projection value based on the linear relationship.

  \begin{asparadesc}
  \item[Preparation:] Same as E1.  Confirm every $N$ in
    \texttt{SCALING\_N\_VALUES} has an entry in
    \texttt{N\_TO\_T}.

  \item[Execution:]
    \texttt{uv run python run\_experiments.py --scaling-only}.

  \item[Results:] Inspect
    \texttt{scaling\_results.png} and
    \texttt{scaling\_results.json}.  Pass criteria: fidelity $\geq 0.85$ across all $N$, and the projected TLP
    \texttt{PSolve} time at $T(N)$ exceeds the HHL time at every
    $N$ (C3).  Absolute times vary with the host CPU.
  \end{asparadesc}
\end{compactitem}

\subsection{Notes on Reusability}

The protocol is not HHL-specific, HHL is one example quantum
algorithm.  \texttt{PythonGarbledTLP} generates a TLP for any
quantum circuit that produces a classical output, so the same
construction works for HHL, VQE, QAOA, Shor's, Grover's, etc. Use cases include verifiability of delegated
quantum computation, fair release of a quantum-computation result
on a cloud backend and so on.

Two parameters can be altered for different use cases:

\begin{compactitem}
\item \textbf{TLP time parameter $t$ (or $T$).}  Set to the value
  that makes \texttt{PSolve} just exceed the
  quantum-computation time on the target hardware.  Use
  \texttt{--tlp-only} to measure the per-unit-$t$ slope on a new
  machine or with a new wrapped circuit, and
  \texttt{--scaling-only} to project that slope to large $t$ via
  the linear fit (the linearity is intrinsic to the TLP and
  holds for any wrapped circuit, not just HHL).
\item \textbf{Quantum problem instance.}  For HHL this is the
  matrix $A$ and right-hand side $b$; set \texttt{HHL\_A} and
  \texttt{HHL\_B} at the top of \texttt{run\_experiments.py}
  (the system is auto-Hermitianized by
  \texttt{hermitianize\_matrix}).  For other algorithms,
  replace the call inside \texttt{run\_scaling\_experiment}
  with the algorithm's own estimate/measurement routine.
\end{compactitem}

\noindent Further optional customizations: swap the TLP's
sequential function for different security models
(SHA-256 default, lattice-based in
\texttt{timelock\_puzzle/}\-\texttt{lattice\_sequential\_function.py},
POSW in
\texttt{timelock\_puzzle/}\-\texttt{posw\_lattice.py}; custom
functions implement the
\texttt{SequentialFunction} interface in
\texttt{timelock\_puzzle/}\-\texttt{sequential\_functions.py}),
and substitute the dependencies (HHL or CRGC libraries) by editing
the imports at the top of \texttt{run\_experiments.py}.

\subsection{Version}
Based on the LaTeX template for Artifact Evaluation V20220926.

%%%%%%%%%%%%%%%%%%%%%%%%%%%%%%%%%%%%%%%%%%%%%%%%%%%%%%%%%%%%%%%%%%%%%

%% This template is based on CCS 2025 template. The AEC Chair for CCS 2026 was:
%% Jérémie Decouchant, Delft University of Technology
%% Fengwei Zhang, Southern University of Science and Technology
%% Jun Dai, Worcester Polytechnic Institute

%% This template was created by David Balenson and Ning Zhang for the ACM CCS 2025
%%
%% The AEC Chair for CCS 2025 was:
%% David Balenson (USC Information Science Institute)
%% Ning Zhang (Washington University in St. Louis)

%% This template is based on CCS 2024 template, 
%%  https://www.sigsac.org/ccs/CCS2024/assets/appendix/ccs2024-ae.zip
%%  the AEC chair for CCS 2024 was 
%%  Eric Eide, University of Utah
%% the template was originally based on the "USENIX Security 2023 AE Template."
%% https://secartifacts.github.io/usenixsec2023/appendix/usesec23-ae-latex.zip
%%
%% The AEC Co-chairs for USENIX Security 2023 were:
%% Cristiano Giuffrida (Vrije Universiteit Amsterdam)
%% Anjo Vahldiek-Oberwagner (Intel Labs)

%% The USENIX Security AE template was based on the
%% % Artifact Appendix Template for EuroSys'22 AE
%% https://sysartifacts.github.io/eurosys2022/appendix/EuroSys22_ArtifactAppendix_template.tex
%%
%% The AEC Co-chairs for EuroSys 2022 were:
%% Natacha Crooks (UC Berkeley)
%% Solal Pirelli (EPFL)
%% Salvatore Signorello (University of Lisbon)
%% Anjo Vahldiek-Oberwagner (Intel Labs)

%% The above were based on:

%% LaTeX template for Artifact Evaluation V20220119
%%
%% Original Authors
%% * Grigori Fursin (cTuning foundation, France) 2014-2020
%% * Bruce Childers (University of Pittsburgh, USA) 2014
%%
%% Modified by
%% * Clementine Maurice (CNRS, France) 2021-2022
%% * Cristiano Giuffrida (Vrije Universiteit Amsterdam, Netherlands) 2021-2022
%%
%% (C) Copyright 2014-2022
%%
%% CC BY 4.0 license
%%

\end{document}
\endinput
